\textbf{6.\ \textsc{Pull-up bar rope} (5 points)} --- \textit{Jaan Kalda.}
A rope is thrown over a horizontal pull-up bar at height $H$ above the
frictionless floor. One end of the rope is tied to a heavy weight resting on
the floor; the other end is pulled so that the length of the rope between the
bar and the weight decreases at a constant rate $v$. At a given moment, this
segment makes an angle $\alpha$ with the vertical (see figure); the weight is
still in contact with the floor. Gravitational acceleration is $g$.

\begin{center}\includegraphics[width=0.5\linewidth]{figures/NBPhO_2026_Q6_fig1.png}\end{center}

\textbf{i)} \textit{(1 point)} Find the speed $u$ of the weight.

\textbf{ii)} \textit{(1 point)} Find the magnitude of the acceleration of the
weight.

\textbf{iii)} \textit{(3 points)} Find the angle $\alpha_0$ at which the
weight lifts off the floor.
